%% ma: L10-B12-0054
%% lop: 10
%% chuong: 5
%% bai: 12
%% dang: 10.12.5
%% loai: TN4
%% muc: VD
%% dap_an: "A"
%% nguon: "chuyen-de-cac-so-dac-trung-cua-mau-so-lieu-khong-ghep-nhom-toan-10-chuong-trinh-moi.pdf (GV. Phan Nhật Linh) (P01), Bài 01 Phần I Câu 38"
%% kien_thuc: "Lớp 10 Bài 12 (quy tròn số), đổi đơn vị, kí hiệu khoa học"
%% trang_thai: chinh-thuc
%% ghi_chu: "Tự duyệt theo yêu cầu giáo viên 22/10/2026."
\begin{ex}
Các nhà khoa học đang nghiên cứu một loại máy bay có thể bay với tốc độ gấp bảy lần tốc độ ánh sáng. Biết tốc độ ánh sáng là $299\,792\,458$ m/s. Nếu tốc độ ánh sáng được làm tròn đến hàng nghìn (đơn vị km/s) thì trong một ngày ($24$ giờ) máy bay đó bay được bao nhiêu km?
\choice{\True $1{,}8144\cdot10^{11}$}{$2{,}592\cdot10^{10}$}{$1{,}8131\cdot10^{11}$}{$5{,}04\cdot10^{7}$}
\loigiai{Một ngày có $24\cdot60\cdot60=86\,400$ giây.

$299\,792\,458$ m/s $=299\,792{,}458$ km/s; làm tròn đến hàng nghìn: $300\,000$ km/s.

Quãng đường: $7\cdot300\,000\cdot86\,400=1{,}8144\cdot10^{11}$ (km). Chọn A.}
\end{ex}
